=======
inverse problems
=======

## Preconditioners for krylov solvers for the H p = g subproblem in pde constrained optimization 
https://repositories.lib.utexas.edu/server/api/core/bitstreams/95072b2e-e489-4026-b5b9-11ab4e12fdd7/content
says Regularisation preconditioning is the most common
but this is only good for when a small number of parameters effect the data misfit term because the operator is "compact" 
I think that this is actually PERFECT for LS methods - since we only really have dependance on the boundary DOFs -- so the regularisation preconditioning should be good....
I am not sure yet if we actually have to add the regularisation term to the objective though
(page 68)

==========
Level sets
==========

## allaires students thesis (also involves that second order allaire paprer)
"Second-order derivatives for shape optimization with a level-set method"
\cite{vie2016second}
very interesting and relevatn
does a cg-newton method and a full newton method for compliance probklems 
Now, they show sorta shit results for discretise then optimise and a biut better for optimise then discretise 
but they say that there discretise then optimise shit results could be due to the fact that they use ersatz 
we could use cutfem..... 
also, they do not play around with the forcing criteria for the cg iterates - i think therefore it is not that approprait
the cg should lie somehere btwn newton and gradient
i think the only reason it would is becaus eof the trust region steps.. 
lets do a bit more investigation down this line... .

\cite{Allaire2016}
the article (not thesis) for normal boundary perturbations

## NGSolve second order stuff 
Full newton for shape opt - BOUNDARY FITTED MESH
\cite{Gangl2020}
shows much faster convergence for H1 regularsised second order for a thermal problem (data misfit objective)
(unconstrained) ? 

\cite{Etling2018}
Similar to above - boundary fitted & onoy very simple example

Homotopy newton shape opt
\cite{Gangl2024}

### Implicit differentiation with second-order derivatives and benchmarks in finite-element-based differentiable physics
shows a very simple shape opt to have benifit with newton - these are the JAXFEM guys

\cite{Hghj2025}
Maute paper - not second order but density transitioning to level sets - i think this would be the right way to introduce second order to LS methods... 



=====
SIMP
=====

## Computing multiple solutions in simp with a second order method ---use for talking about continuation...
\cite{Papadopoulos2021} 
They use the method we are proposing with simp -- and do something to the newton iterations to get multiple solutions i belive - not hessian free though i think 

## SQP simp 
\cite{RojasLabanda2016}
Shows reduced iterations when using a newton like SQP method (full (modifired) hessian built - obviously expensice )

## SQP simp same as above - ole sigmund
\cite{RojasLabanda2017}
shows the above method requiring fewer iterations than MMA 

## another simp sqp but a bit pooror quality papoer 
\cite{Liao2022}
Shows slightly improved convergence 

## solving simp newton on a reduced space so solves are cheaper 
\cite{Evgrafov2014}

=====